% TOP_PROBLEM: {"id": 4800032, "title": "Multiple ergodic averages — Problem 32", "category_id": 11, "category": {"id": 11, "name": "dynamical_systems", "display_name": "Dynamical Systems", "description": "Problems about long-term behavior of deterministic systems, Hamiltonian dynamics, and periodic orbits.", "slug": "dynamical-systems", "order_index": 11, "created_at": "2026-07-31T15:26:25.671Z"}, "status": "partially_solved"}
% FIRST_POSTED: null
% ENTRY_KIND: statement_only
% Imported from: batch15.tex; UnsolvedMath ID 4800032
\documentclass[11pt]{article}
\usepackage[margin=25mm]{geometry}
\usepackage{fontspec}
\setmainfont{Latin Modern Roman}
\usepackage{amsmath,amssymb,mathtools}
\usepackage{unicode-math}
\setmathfont{Latin Modern Math}
\usepackage{xurl,hyperref}
\hypersetup{colorlinks=true,urlcolor=blue}
\setcounter{secnumdepth}{0}
\setlength{\parindent}{0pt}
\setlength{\parskip}{5pt}
\setlength{\emergencystretch}{3em}
\newcommand{\sref}[2]{\hyperlink{source:#1:#2}{[S#2]}}
\newcommand{\cataloguescope}[1]{\paragraph{Recorded target.}#1}
\begin{document}
% UnsolvedMath ID 4800032; source catalogue code AMR-047-0032
\setcounter{section}{4800031}
\section{Multiple ergodic averages — Problem 32}\label{4800032}
{\small\sffamily UnsolvedMath ID 4800032 · Dynamical Systems}
\subsection{Definitions and mathematical statement}

\paragraph{Shared notation.}$\mathbb N=\{1,2,\ldots\}$, and $\mathbb Z,\mathbb Q,\mathbb R,\mathbb C$ denote the integers, rationals, reals and complex numbers. Any other conventions are specified locally.


A system is a probability space $(X,\mathcal X,\mu)$ with invertible measure-preserving transformations; multiple transformations are assumed to commute. Write $T^j f=f\circ T^j$, and $T^j A$ for the image of a measurable set under $T^j$. A system with one transformation is ergodic if every invariant measurable set has measure zero or one. The space $L^\infty(\mu)$ consists of essentially bounded complex measurable functions. Let $(\sigma_j)_{j\ge1}$ be a nonincreasing sequence in $(0,1]$ with $j\sigma_j\to\infty$. On an auxiliary probability space $(\Omega,\mathcal F,\mathbb P)$ take independent Bernoulli variables $X_j$ with $\mathbb P(X_j=1)=\sigma_j$. Enumerate the selected positive integers increasingly: $a_n(\omega)=\min\{k:\sum_{j=1}^kX_j(\omega)=n\}$. This sequence is infinite almost surely. For a transformation $T$, let $\mathcal I_T$ be its invariant $\sigma$-algebra. Conditional expectation $\mathbb E(f\mid\mathcal I_T)$ is the orthogonal projection in $L^2$ onto the invariant functions.
\paragraph{Statement recorded in the supplied source.}
Is there a single event $\Omega_0$ of probability one such that for every $\omega\in\Omega_0$, every commuting system $(X,\mathcal X,\mu,T,S)$ and all bounded $f,g$, \[\lim_{N\to\infty}\frac1N\sum_{n=1}^NT^nf\,S^{a_n(\omega)}g=\mathbb E(f\mid\mathcal I_T)\mathbb E(g\mid\mathcal I_S)\quad\text{in }L^2(\mu)?\] Moreover, when $\sigma_j=j^{-\alpha}$ for any $\alpha\in(0,1)$, does a universal probability-one event give the same convergence for $\mu$-almost every point of every such system? \sref{4800032}{2}
The first question is reproduced with the stated nonincreasing scope. That literal formulation has a boundary obstruction: if $\sigma_j=1$ for every $j$, then $a_n=n$. On the probability circle $\mathbb R/\mathbb Z$ with Lebesgue measure, let $T=S$ be an irrational rotation, take $f(x)=e^{2\pi i x}$ and $g=\bar f$. This is an ergodic system, so both conditional expectations on the right vanish, whereas $T^nf\,S^ng=1$ for every $n$. Thus this literal first formulation has a negative answer; the separate power-law question retains its stated scope.
\subsection{Short English statement}
Do nonincreasing random-selection probabilities with $n\sigma_n\to\infty$ give the universal bilinear mean limit, and do all power-law selections also give the pointwise limit?
\subsection{Sources}
\begin{description}
\item[\hypertarget{source:4800032:1}{S1}] ULAM.AI and the UnsolvedMath Contributors, \emph{UnsolvedMath: A Curated Collection of Open Mathematics Problems}, record 4800032 (AMR-047-0032). CC BY 4.0. \url{https://huggingface.co/datasets/ulamai/UnsolvedMath}.
\item[\hypertarget{source:4800032:2}{S2}] N. Frantzikinakis, Some open problems on multiple ergodic averages (2016), Section 7 complete random-sequence setup, universal-almost-sure definition and full Problem 32, pp.34--35. \url{https://arxiv.org/pdf/1103.3808v3}.
\end{description}
\label{end:4800032}
\end{document}
